\documentclass[12pt,a4paper]{article}
\usepackage{amsmath,amssymb,amsthm,hyperref,geometry}
\usepackage{enumitem,booktabs,array}
\usepackage[ngerman]{babel}
\geometry{margin=2.5cm}

% -------------------------------------------------------
% Theorem-Umgebungen (Deutsch)
% -------------------------------------------------------
\newtheorem{theorem}{Satz}[section]
\newtheorem{lemma}[theorem]{Lemma}
\newtheorem{korollar}[theorem]{Korollar}
\newtheorem{proposition}[theorem]{Proposition}
\newtheorem{vermutung}[theorem]{Vermutung}
\newtheorem{definition}[theorem]{Definition}

\theoremstyle{remark}
\newtheorem{bemerkung}[theorem]{Bemerkung}
\newtheorem{beispiel}[theorem]{Beispiel}

% -------------------------------------------------------
% Metadaten
% -------------------------------------------------------
\title{\textbf{Die Riemann-Hypothese: Ein analytischer Ansatz}\\[0.5em]
\large Nullstellenfreie Gebiete, Explizite Formeln und die kritische Gerade}
\author{Michael Fuhrmann}
\date{11.\ März 2026}

\hypersetup{
  pdftitle={Die Riemann-Hypothese: Ein analytischer Ansatz},
  pdfauthor={Michael Fuhrmann},
  colorlinks=true,
  linkcolor=blue,
  citecolor=blue
}

\begin{document}
\maketitle
\begin{abstract}
Die Riemann-Hypothese (RH), von Bernhard Riemann 1859 formuliert, behauptet, dass alle
nichttrivialen Nullstellen der Riemann-Zeta-Funktion $\zeta(s)$ auf der kritischen Geraden
$\operatorname{Re}(s)=\tfrac{1}{2}$ liegen.  Trotz mehr als 160~Jahren intensiver Forschung
und der Verifikation von über $10^{13}$ Nullstellen blieb die Vermutung unbewiesen und
gehört zu den Millennium-Problemen des Clay Mathematics Institute.

Dieses Papier gibt einen Überblick über die wichtigsten analytischen Werkzeuge:
nullstellenfreie Gebiete von de~la~Vallée~Poussin (1899) und die
Korobov--Vinogradov-Methode, die explizite Formel für $\psi(x)$, Hardy--Littlewoo\-ds
Ergebnisse über Nullstellen auf der kritischen Geraden, numerische Verifikationen von
Odlyzko und van~de~Lune sowie die Verallgemeinerte Riemann-Hypothese (GRH) für
Dirichlet-$L$-Funktionen.
\end{abstract}

\tableofcontents
\newpage

% ================================================================
\section{Historischer Hintergrund}
% ================================================================

In seiner berühmten Abhandlung \textit{Über die Anzahl der Primzahlen unter einer gegebenen
Grösse} von 1859 führte Riemann die analytische Fortsetzung der Euler--Dirichlet-Reihe
\[
  \zeta(s) \;=\; \sum_{n=1}^{\infty} \frac{1}{n^s}, \qquad \operatorname{Re}(s) > 1,
\]
zu einer meromorphen Funktion auf ganz $\mathbb{C}$ (mit einem einfachen Pol bei $s=1$) ein.
Riemann vermutete – ohne Beweis –, dass alle nichttrivialen Nullstellen auf der Geraden
$\operatorname{Re}(s)=\tfrac{1}{2}$ liegen.

Diese Abhandlung legte den Grundstein für den Primzahlsatz (PNT), der 1896 unabhängig von
Hadamard und de~la~Vallée~Poussin bewiesen wurde.  Das Clay~Institute erklärte die RH im
Jahr 2000 zu einem der sieben Millennium-Probleme (Preisgeld: 1\,000\,000 USD).

% ================================================================
\section{Die Riemann-Zeta-Funktion}
% ================================================================

\subsection{Analytische Fortsetzung und Funktionalgleichung}

Für $\operatorname{Re}(s)>1$ gilt das Euler-Produkt
\[
  \zeta(s) \;=\; \prod_{p\;\text{prim}} \frac{1}{1 - p^{-s}},
\]
das den Fundamentalsatz der Arithmetik codiert.  Die analytische Fortsetzung liefert eine
auf $\mathbb{C}\setminus\{1\}$ holomorphe Funktion mit der Funktionalgleichung
\begin{equation}\label{eq:func}
  \xi(s) \;=\; \xi(1-s),
  \qquad
  \xi(s) \;:=\; \tfrac{1}{2}\,s(s-1)\,\pi^{-s/2}\,\Gamma\!\left(\tfrac{s}{2}\right)\zeta(s).
\end{equation}
Die vervollständigte $\xi$-Funktion ist ganz, auf der reellen Achse und auf der kritischen
Geraden reellwertig sowie symmetrisch: $\xi(s)=\xi(\bar{s})$.

\subsection{Triviale und nichttriviale Nullstellen}

Die Funktionalgleichung erzwingt \emph{triviale Nullstellen} bei $s=-2,-4,-6,\ldots$
(durch die Pole von $\Gamma(s/2)$).  Alle weiteren Nullstellen sind \emph{nichttrivial}
und liegen im \emph{kritischen Streifen} $0<\operatorname{Re}(s)<1$.

\begin{vermutung}[Riemann-Hypothese, 1859]
  Jede nichttriviale Nullstelle $\rho$ von $\zeta(s)$ erfüllt $\operatorname{Re}(\rho)=\tfrac{1}{2}$.
\end{vermutung}

Nichttriviale Nullstellen treten in konjugierten Paaren auf:
\[
  \zeta(\rho)=0 \;\Longrightarrow\; \zeta(\bar\rho)=0=\zeta(1-\rho)=\zeta(1-\bar\rho).
\]
Man schreibt sie üblicherweise als $\rho=\tfrac{1}{2}+i\gamma$; die ersten Imaginärteile
sind $\gamma_1\approx14{,}135$, $\gamma_2\approx21{,}022$, $\gamma_3\approx25{,}011$.

% ================================================================
\section{Nullstellenfreie Gebiete}
% ================================================================

\subsection{De la Vallée Poussin (1899)}

\begin{theorem}[de~la~Vallée~Poussin, 1899]\label{thm:dvp}
  Es existiert eine absolute Konstante $c>0$, sodass $\zeta(s)\neq0$ für
  \[
    \sigma \;:=\; \operatorname{Re}(s) \;\geq\; 1 - \frac{c}{\log(|t|+2)},
    \qquad t := \operatorname{Im}(s).
  \]
\end{theorem}

\begin{proof}[Beweisskizze]
  Die trigonometrische Ungleichung
  \[
    3 + 4\cos\theta + \cos(2\theta) = 2(1+\cos\theta)^2 \;\geq\; 0
  \]
  ergibt für $\theta=t\log p$ und $\sigma>1$:
  \[
    3\log\zeta(\sigma) + 4\,\operatorname{Re}\!\left[-\frac{\zeta'}{\zeta}(\sigma+it)\right]
    + \operatorname{Re}\!\left[-\frac{\zeta'}{\zeta}(\sigma+2it)\right] \;\geq\; 0.
  \]
  Wenn $\zeta(\rho_0)=0$ mit $\rho_0=\beta+it_0$, liefert der Pol-Residuen-Vergleich
  beim Grenzübergang $\sigma\to1^+$ die Schranke $\beta\leq1-c/\log t_0$.
\end{proof}

Dieses nullstellenfreie Gebiet war das entscheidende Werkzeug der Primzahlsatz-Beweise von
1896.

\subsection{Korobov--Vinogradov-Methode (1958)}

\begin{theorem}[Korobov~\cite{Korobov1958}; Vinogradov~\cite{Vinogradov1958}]
  Es gibt $c>0$ mit $\zeta(s)\neq0$ für
  \begin{equation}\label{eq:kv}
    \sigma \;\geq\; 1 - \frac{c}{(\log|t|)^{2/3}(\log\log|t|)^{1/3}},
    \qquad |t| \geq 3.
  \end{equation}
\end{theorem}

Die Verbesserung von $(\log t)^{-1}$ auf $(\log t)^{-2/3}(\log\log t)^{-1/3}$ beruht auf
der Schätzung der Exponentialsumme $\sum_{n\leq N}n^{-\sigma-it}$ mit der
Weyl--van~der~Corput--Vinogradov-Methode.

\begin{bemerkung}
  Gleichung~\eqref{eq:kv} ist das beste bekannte unbedingte nullstellenfreie Gebiet.
  Unter der RH erstreckt es sich auf die gesamte Halbebene $\sigma>\tfrac{1}{2}$.
\end{bemerkung}

\subsection{Folgen für den Primzahlsatz}

Das nullstellenfreie Gebiet steuert den Fehlerterm im PNT direkt:
\[
  \pi(x) \;=\; \operatorname{Li}(x) + O\!\left(x\exp\!\left(-c\,\frac{(\log x)^{3/5}}{(\log\log x)^{1/5}}\right)\right).
\]
Die RH würde den optimalen Fehlerterm $O(\sqrt{x}\log x)$ liefern.

% ================================================================
\section{Explizite Formeln}
% ================================================================

\subsection{Von Mangoldts explizite Formel}

Sei $\psi(x)=\sum_{n\leq x}\Lambda(n)$ die Chebyshev-Funktion mit der von-Mangoldt-Funktion
$\Lambda$ ($\Lambda(p^k)=\log p$, sonst null).

\begin{theorem}[Explizite Formel, von~Mangoldt 1895]
  Für $x>1$, $x$ keine Primzahlpotenz:
  \begin{equation}\label{eq:explicit}
    \psi(x) \;=\; x \;-\; \sum_{\rho}\frac{x^{\rho}}{\rho} \;-\; \log(2\pi)
              \;-\; \tfrac{1}{2}\log\!\left(1 - x^{-2}\right),
  \end{equation}
  wobei die Summe über alle nichttrivialen Nullstellen $\rho$ von $\zeta$, geordnet nach
  $|\operatorname{Im}(\rho)|$, läuft.
\end{theorem}

\begin{bemerkung}
  Gilt die RH, so haben alle $x^\rho=x^{1/2}e^{i\gamma\log x}$ den Betrag $\sqrt{x}$,
  und man erhält $\psi(x)=x+O(\sqrt{x}\log^2 x)$.  Unbedingt ist der beste bekannte
  Fehlerterm um einen Logarithmusfaktor größer.
\end{bemerkung}

\subsection{Riemanns Nullstellenzählformel}

Die Anzahl der Nullstellen mit $0<\operatorname{Im}(\rho)\leq T$ beträgt
\begin{equation}\label{eq:counting}
  N(T) \;=\; \frac{T}{2\pi}\log\frac{T}{2\pi} - \frac{T}{2\pi} + O(\log T),
\end{equation}
bekannt als Riemann--von~Mangoldt-Formel.  Die Nullstellen werden mit wachsendem $T$
immer dichter.

% ================================================================
\section{Nullstellen auf der kritischen Geraden}
% ================================================================

\subsection{Hardys Satz (1914)}

\begin{theorem}[Hardy, 1914~\cite{Hardy1914}]
  $\zeta(\tfrac{1}{2}+it)$ besitzt unendlich viele reelle Nullstellen, d.\,h.\
  $\zeta$ hat unendlich viele Nullstellen auf der kritischen Geraden.
\end{theorem}

Hardys Beweis nutzte das Integral $\int_0^\infty\Xi(t)\cos(t\log u)/t\,dt$ mit
$\Xi(t)=\xi(\tfrac{1}{2}+it)\in\mathbb{R}$ und zeigte unendlich viele Vorzeichenwechsel.

\subsection{Hardy--Littlewood (1921)}

\begin{theorem}[Hardy--Littlewood, 1921~\cite{HardyLittlewood1921}]
  Die Anzahl $N_0(T)$ der Nullstellen von $\zeta(\tfrac{1}{2}+it)$ mit $0\leq t\leq T$
  erfüllt $N_0(T)\geq AT\log T$ für eine absolute Konstante $A>0$.
\end{theorem}

\subsection{Levinsons Methode (1974)}

\begin{theorem}[Levinson, 1974~\cite{Levinson1974}]
  Mindestens $\tfrac{1}{3}$ aller nichttrivialen Nullstellen von $\zeta(s)$ liegen auf der
  kritischen Geraden.
\end{theorem}

Levinsons Mollifikationsmethode multipliziert $\zeta(\tfrac{1}{2}+it)$ mit einem
Dirichlet-Polynom $M(\tfrac{1}{2}+it)$, das die Funktion „glättet" und
Vorzeichenwechsel leichter sichtbar macht.

\subsection{Conrey und spätere Verbesserungen}

\begin{theorem}[Conrey, 1989~\cite{Conrey1989}]
  Mindestens $\tfrac{2}{5}$ aller nichttrivialen Nullstellen liegen auf
  $\operatorname{Re}(s)=\tfrac{1}{2}$.
\end{theorem}

Neuere Arbeiten mit höheren Mollifiern und verschobenen Momenten brachten diesen Anteil
auf über $41\%$, weit vom erforderlichen $100\%$ der RH entfernt.

% ================================================================
\section{Numerische Verifikation}
% ================================================================

\subsection{Riemann--Siegel-Formel}

Für praktische Berechnungen verwendet man die reellwertige Funktion
\[
  Z(t) \;=\; e^{i\vartheta(t)}\zeta\!\left(\tfrac{1}{2}+it\right) \;\in\;\mathbb{R},
  \quad
  \vartheta(t) = \operatorname{Im}\log\Gamma\!\left(\tfrac{1}{4}+\tfrac{it}{2}\right)
                 - \tfrac{t}{2}\log\pi.
\]
Vorzeichenwechsel von $Z(t)$ liefern Nullstellen von $\zeta$ auf der kritischen Geraden.

\subsection{Odlyzko und van de Lune}

\begin{theorem}[van~de~Lune, te~Riele, Winter, 1986~\cite{vanDeLune1986}]
  Alle $1{,}5\times10^9$ Nullstellen mit $0<t\leq T_0$ sind einfach und liegen auf der
  kritischen Geraden.
\end{theorem}

\begin{theorem}[Platt--Trudgian, 2021~\cite{PlattTrudgian2021}]
  Die ersten $3{,}0000001\times10^{12}$ nichttrivialen Nullstellen von $\zeta$ sind
  alle einfach und liegen auf $\operatorname{Re}(s)=\tfrac{1}{2}$.
\end{theorem}

\begin{bemerkung}
  Numerische Verifikation, so umfangreich sie auch ist, kann die RH nicht beweisen:
  Man kann eine Nullstelle außerhalb der kritischen Geraden bei astronomisch großer Höhe
  nicht ausschließen.  Die statistischen Eigenschaften der Nullstellen bei
  $T\approx10^{22}$ stimmen verblüffend gut mit der GUE-Zufallsmatrixstatistik überein
  (Montgomery--Odlyzko-Gesetz) und liefern starke heuristische Evidenz.
\end{bemerkung}

% ================================================================
\section{Die Verallgemeinerte Riemann-Hypothese}
% ================================================================

\subsection{Dirichlet-$L$-Funktionen}

Sei $\chi$ ein Dirichlet-Charakter modulo $q$.  Die zugehörige $L$-Funktion
\[
  L(s,\chi) \;=\; \sum_{n=1}^{\infty} \frac{\chi(n)}{n^s} \;=\; \prod_p \frac{1}{1-\chi(p)p^{-s}}
\]
setzt sich analytisch auf ganz $\mathbb{C}$ fort (ganz falls $\chi\neq\chi_0$).

\begin{vermutung}[Verallgemeinerte Riemann-Hypothese (GRH)]
  Für jeden primitiven Dirichlet-Charakter $\chi$ liegen alle nichttrivialen Nullstellen
  von $L(s,\chi)$ auf $\operatorname{Re}(s)=\tfrac{1}{2}$.
\end{vermutung}

\subsection{Konsequenzen der GRH}

Die GRH hat weitreichende arithmetische Folgen:
\begin{itemize}[nosep]
  \item \textbf{Primzahlen in arithmetischen Progressionen:}
    $\pi(x;q,a) = \operatorname{Li}(x)/\phi(q) + O(\sqrt{x}\log(qx))$.
  \item \textbf{Artins Vermutung:} Jede nicht-quadratische ganze Zahl $a\neq\pm1$
    ist primitive Wurzel modulo unendlich vieler Primzahlen (Hooley 1967, bed.\ GRH).
  \item \textbf{Millers Primzahltest:} Deterministisch polynomialzeit (Miller 1976,
    bed.\ GRH).
  \item \textbf{Klassenzahlen:} Gaußsche Klassenzahlvermutung $h(-d)\to\infty$
    (Goldfeld--Gross--Zagier, teilweise bed.).
  \item \textbf{Goldbach:} Jede ungerade Zahl $n>5$ ist Summe dreier Primzahlen
    (Helfgott 2013 unbedingt; GRH liefert drastisch kürzeren Beweis).
\end{itemize}

\subsection{Große Riemann-Hypothese}

Die Große Riemann-Hypothese verallgemeinert die GRH auf alle automorphen $L$-Funktionen
im Langlands-Programm, einschließlich symmetrischer Potenzen $L(s,\operatorname{sym}^k f)$
holomorpher Neuformen und $L$-Funktionen von Motiven.

% ================================================================
\section{Verbindungen zur Zufallsmatrixtheorie}
% ================================================================

Montgomerys Paarkorrelations-Vermutung~\cite{Montgomery1973} besagt, dass die normierten
Nullstellenabstände $\tilde\gamma_n = \gamma_n(\log\gamma_n)/(2\pi)$ die
Paarkorrelationsdichte
\[
  1 - \left(\frac{\sin\pi u}{\pi u}\right)^2
\]
besitzen – identisch mit der GUE-Eigenwert-Paarkorrelation.  Odlyzkos numerische Experimente
bestätigten dies mit bemerkenswerter Präzision.  Existierte ein selbstadjungierter Operator
(Hilbert--Pólya-Vermutung), dessen Eigenwerte die Imaginärteile der Nullstellen sind, so
würde die RH unmittelbar folgen.

% ================================================================
\section{Aktueller Stand und Perspektiven}
% ================================================================

\begin{itemize}[nosep]
  \item \textbf{Unbedingt:} Nullstellenfreies Gebiet vom Korobov--Vinogradov-Typ;
    unendlich viele Nullstellen auf der kritischen Geraden; mindestens $41\%$ dort.
  \item \textbf{Numerisch:} Über $10^{13}$ Nullstellen auf der kritischen Geraden verifiziert.
  \item \textbf{Bed.\ RH:} Optimaler PNT-Fehlerterm, vollständige Artin-Vermutung,
    effektive Siegel-Nullstellen-Schranken.
  \item \textbf{Offen:} Kein Ansatz bekannt, der auch nur eine einzige Nullstelle
    außerhalb der kritischen Geraden ausschließt.
  \item \textbf{Spektralansatz:} Berry--Keating-Vermutung $H=xp$ (Quantisierung des
    klassischen Hamiltonians $H=xp$ im Phasenraum).
  \item \textbf{Absolute Geometrie:} Weilsche Vermutungen (Deligne 1974 für Funktionenkörper)
    und der Körper $\mathbb{F}_1$ als möglicher Weg.
\end{itemize}

% ================================================================
\section{Schlussfolgerung}
% ================================================================

Die Riemann-Hypothese bleibt eines der tiefsten ungelösten Probleme der Mathematik.
Die analytische Theorie von $\zeta(s)$ ist reich entwickelt: Nullstellenfreie Gebiete,
explizite Formeln und Verbindungen zur Zufallsmatrixtheorie deuten alle auf die
Wahrheit der RH hin, und die numerische Evidenz ist überwältigend.  Dennoch scheint
ein Beweis fundamental neue Ideen zu erfordern – ob aus der Spektraltheorie, der
algebraischen Geometrie oder einer völlig unerwarteten Richtung.

\begin{thebibliography}{99}

\bibitem{Riemann1859}
B.~Riemann,
\textit{Über die Anzahl der Primzahlen unter einer gegebenen Grösse},
Monatsberichte der Berliner Akademie (1859).

\bibitem{Hadamard1896}
J.~Hadamard,
\textit{Sur la distribution des zéros de la fonction $\zeta(s)$},
Bull.\ Soc.\ Math.\ France \textbf{24} (1896), 199--220.

\bibitem{delaValleePoussin1899}
C.-J.~de~la~Vallée~Poussin,
\textit{Sur la fonction $\zeta(s)$ de Riemann},
Mém.\ Couronnés Acad.\ Roy.\ Belg.\ \textbf{59} (1899).

\bibitem{Korobov1958}
N.~M.~Korobov,
\textit{Schätzungen trigonometrischer Summen und ihre Anwendungen},
Uspekhi Mat.\ Nauk \textbf{13} (1958), 185--192.

\bibitem{Vinogradov1958}
I.~M.~Vinogradov,
\textit{Eine neue Schätzung von $\zeta(1+it)$},
Izv.\ Akad.\ Nauk SSSR, Ser.\ Mat.\ \textbf{22} (1958), 161--164.

\bibitem{Hardy1914}
G.~H.~Hardy,
\textit{Sur les zéros de la fonction $\zeta(s)$ de Riemann},
C.\ R.\ Acad.\ Sci.\ Paris \textbf{158} (1914), 1012--1014.

\bibitem{HardyLittlewood1921}
G.~H.~Hardy und J.~E.~Littlewood,
\textit{The zeros of Riemann's zeta-function on the critical line},
Math.\ Z.\ \textbf{10} (1921), 283--317.

\bibitem{Levinson1974}
N.~Levinson,
\textit{More than one third of zeros of Riemann's zeta-function are on $\sigma=\tfrac{1}{2}$},
Adv.\ Math.\ \textbf{13} (1974), 383--436.

\bibitem{Conrey1989}
J.~B.~Conrey,
\textit{More than two fifths of the zeros of the Riemann zeta function are on the critical line},
J.\ Reine Angew.\ Math.\ \textbf{399} (1989), 1--26.

\bibitem{vanDeLune1986}
J.~van~de~Lune, H.~J.~J.~te~Riele und D.~T.~Winter,
\textit{On the zeros of the Riemann zeta function in the critical strip, IV},
Math.\ Comp.\ \textbf{46} (1986), 667--681.

\bibitem{Odlyzko2001}
A.~M.~Odlyzko,
\textit{The $10^{22}$-nd zero of the Riemann zeta function},
in \textit{Dynamical, Spectral, and Arithmetic Zeta Functions}, AMS, 2001.

\bibitem{PlattTrudgian2021}
D.~J.~Platt und T.~S.~Trudgian,
\textit{The Riemann hypothesis is true up to $3\cdot10^{12}$},
Bull.\ Lond.\ Math.\ Soc.\ \textbf{53} (2021), 792--797.

\bibitem{Montgomery1973}
H.~L.~Montgomery,
\textit{The pair correlation of zeros of the zeta function},
Proc.\ Sympos.\ Pure Math.\ \textbf{24}, AMS (1973), 181--193.

\bibitem{Titchmarsh1986}
E.~C.~Titchmarsh,
\textit{The Theory of the Riemann Zeta-Function}, 2.~Aufl.,
Oxford University Press, 1986.

\bibitem{Iwaniec2004}
H.~Iwaniec und E.~Kowalski,
\textit{Analytic Number Theory},
AMS Colloquium Publications \textbf{53}, 2004.

\end{thebibliography}

\end{document}
